\chapter{Gaussian Parcels in a Straight Channel}\label{ch:gaussian}

\section{From individual particles to represented populations}
The second stage adds hydrodynamic feedback while retaining a simple
carrier geometry. A parcel represents $N_i$ identical spherical particles
whose centers are approximated by a common planar state. Its mass and
hydrodynamic reaction are weighted by that multiplicity. The Gaussian
distribution describes how the represented force interacts with the
carrier; it does not resolve the surface of a sphere.

The dimensionless fluid domain is $\Omega=[0,3]\times[0,1]$, periodic in
$x$ with stationary no-slip walls at $y=0$ and $y=1$. The channel has no
open inlet or outlet. The initial fluid velocity is
\begin{equation}
 \vect{u}_0(y)=(6y(1-y),0),
 \label{eq:gaussian_initial}
\end{equation}
whose mean streamwise speed is one. A body-force acceleration
$\vect{f}_d=(12/\Rey_H,0)$ maintains this profile in the unladen steady
limit. Parcel initial velocities are zero. The reference depth ratio is
$b=1$, and the evaluated model uses $\Rey_H=20$, $R=1000$, six parcels,
and Gaussian width $\sigma=0.14$.

\section{Coupled equations}
The fluid equations use the inertial pressure convention defined in
\cref{eq:gaussian_scales}:
\begin{align}
 \nabla\cdot\vect{u}&=0,\\
 \partial_t\vect{u}+(\vect{u}\cdot\nabla)\vect{u}
 &=-\nabla p_i+\frac{1}{\Rey_H}\nabla^2\vect{u}
   +\vect{S}+\vect{f}_d.
 \label{eq:gaussian_ns}
\end{align}
The parcel equations are
\begin{equation}
 \dot{\vect{q}}_i=\vect{v}_i,\qquad
 \dot{\vect{v}}_i=\frac{\vect{U}_i-\vect{v}_i}{\St_i}
 +\vect{g}_{\mathrm{eff}},\qquad
 \vect{F}_i=\frac{m_i}{\St_i}(\vect{U}_i-\vect{v}_i).
 \label{eq:gaussian_particle}
\end{equation}
Here $\vect{g}_{\mathrm{eff}}$ is a specified effective acceleration,
including any buoyancy adjustment in the physical mapping. The reaction
source contains the hydrodynamic drag only. Depositing parcel gravity again
as fluid drag would count a different force twice.

\section{Normalized Gaussian transfer}
Let $\vect{x}_a$ denote the fixed transfer-quadrature locations and
$\Delta A_a$ their area weights. A periodized raw Gaussian is evaluated
between $\vect{x}_a$ and each parcel center. In the two-dimensional channel,
its structure is
\begin{equation}
 G_{ai}=\sum_{k\in\mathbb{Z}}
 \exp\left[-\frac{(x_a-q_{ix}+kL_x)^2+(y_a-q_{iy})^2}
 {2\sigma^2}\right].
 \label{eq:raw_kernel}
\end{equation}
The implementation retains a finite image set with distant tails beyond
eight widths after choosing the nearest image. A global Gaussian prefactor
is unnecessary because normalization removes it. Define
\begin{equation}
 Z_i=\sum_aG_{ai}\Delta A_a,\qquad W_{ai}=G_{ai}/Z_i,
 \qquad \sum_aW_{ai}\Delta A_a=1.
 \label{eq:normalized_kernel}
\end{equation}
Both fluid interpolation and force spreading use these weights:
\begin{equation}
 \vect{U}_i=\sum_aW_{ai}\vect{u}_a\Delta A_a,\qquad
 \vect{S}_a=-\sum_iW_{ai}\vect{F}_i.
 \label{eq:shared_transfer}
\end{equation}
Near a wall, $Z_i$ accounts for the actual bounded-domain quadrature.
At an arbitrary collocation point, the numerator uses that point, while
the denominator remains the value from the fixed transfer quadrature.

\subsection{Force and work identities}
Substituting \cref{eq:shared_transfer} into the quadrature sum gives
\begin{equation}
 \sum_a\vect{S}_a\Delta A_a=-\sum_i\vect{F}_i,
 \label{eq:force_identity}
\end{equation}
and
\begin{equation}
 \sum_a\vect{u}_a\cdot\vect{S}_a\Delta A_a
 =-\sum_i\vect{U}_i\cdot\vect{F}_i.
 \label{eq:work_identity}
\end{equation}
The combined interphase contribution to fluid and parcel kinetic-energy
rates is therefore
\begin{equation}
 \mathcal{P}_{\mathrm{drag}}
 =\sum_i(\vect{v}_i-\vect{U}_i)\cdot\vect{F}_i
 =-\sum_i\frac{m_i}{\St_i}|\vect{U}_i-\vect{v}_i|^2\leq0.
 \label{eq:drag_dissipation}
\end{equation}
These are semidiscrete properties of the shared transfer. A trained PINN
with nonzero residual need not satisfy the complete global momentum or
energy balance exactly. External driving can also make total kinetic energy
increase even though the drag contribution is dissipative.

\section{Neural representation and hard constraints}
Two MLPs, each with three hidden layers of 32 tanh units, represent the fluid
correction and parcel histories. The fluid network uses wall-normal position,
time, and three periodic Fourier harmonics in $x$. The ansatz is
\begin{equation}
 \vect{u}(\vect{x},t)=\vect{u}_0(y)
 +\frac{t}{T}y(1-y)\vect{n}_{u}(\vect{x},t).
 \label{eq:fluid_ansatz}
\end{equation}
It enforces the initial velocity and stationary wall velocity. Periodicity
follows from the Fourier features. A time-dependent pressure offset is
removed at $(0,0)$; no independent initial-pressure condition is imposed.

The parcel network uses the relaxation time in its output envelopes:
\begin{align}
 b_i(t)&=1-\exp(-t/\St_i),&
 a_i(t)&=t-\St_i b_i(t),\\
 \vect{q}_i(t)&=\vect{q}_{i0}+t\vect{v}_{i0}
                    +a_i(t)\vect{n}_{q,i}(t),&
 \vect{v}_i(t)&=\vect{v}_{i0}+b_i(t)\vect{n}_{v,i}(t).
 \label{eq:parcel_ansatz}
\end{align}
The envelope imposes the initial states and represents the known response
scale. It supplies no target fluid velocity or target coupled trajectory.
Position and velocity remain separate trainable functions, linked by
$\dot{\vect{q}}_i-\vect{v}_i$ in the objective.

The weighted mean-square residual objective assigns weights 1 to fluid
momentum, 2 to continuity, 5 to parcel kinematics, 1 to parcel acceleration,
and 100 to center confinement. All calculations in this formulation use
float64. \Cref{fig:gaussian_architecture} emphasizes the common quadrature
that couples the representations. The default calculations are
collision-free; independently tested impact primitives are not treated as
a validated collisional suspension model.

\ResultFigure[.96]{08_gaussian_architecture}{Coupled Gaussian formulation.
The two stylized network symbols indicate distinct MLPs, each with three
32-unit hidden layers; their radial arrangement is illustrative and does
not specify literal layer connectivity. Shared quadrature transfers sampled
fluid velocity to parcels and hydrodynamic reaction to the fluid. The
complete objective includes fluid and parcel residuals.}
{fig:gaussian_architecture}

\section{Reduced Stokes-response verification}
The first reduced test, V01, solves the constant-forcing relaxation problem
with dimensionless response time one:
\begin{equation}
 \dot v=1-v,\qquad \dot z=v,\qquad v(0)=z(0)=0.
 \label{eq:stokes_test}
\end{equation}
Its analytical response is
\begin{equation}
 v=1-e^{-t},\qquad z=t-1+e^{-t}.
 \label{eq:stokes_exact}
\end{equation}
The reference arrays are excluded from training and used only after the
physics-constrained fit. The relative velocity and displacement errors are
\VOneVelocity\% and \VOneDistance\%, with independent ODE residual RMS
approximately $1.35\times10^{-3}$. Both the stated relative-error and
residual criteria are satisfied. \Cref{fig:stokes_response} shows the
stored predictions against the analytical response.

\ResultFigure{09_stokes_response}{V01 reduced Stokes response. The PINN
is fitted to the ODE residual and initial conditions; analytical values are
held out. Open markers are subsampled stored predictions for legibility;
the metrics use every evaluation sample. This test verifies the reduced
response equation, not the full two-way channel.}{fig:stokes_response}

\section{Two-way momentum-exchange verification}
The second reduced problem, V02, considers spatially uniform fluid and
particle velocities, fluid mass one, parcel mass loading $\beta=0.25$,
and response time $\St=0.2$. With $u(0)=0$ and $v(0)=1$,
\begin{equation}
 \dot u=-\frac{\beta}{\St}(u-v),\qquad
 \dot v=\frac{u-v}{\St}.
 \label{eq:two_way_test}
\end{equation}
The solution has constant total momentum $u+\beta v=0.25$ and decaying
slip. The relative fluid- and particle-velocity errors are
\VTwoFluid\% and \VTwoParticle\%, respectively. The independent residual
RMS is approximately $6.12\times10^{-3}$.

In this reduced PINN, the output relation $v=1-4u$ imposes total momentum
exactly. Its zero momentum error is consequently a hard-constraint check,
not an independent demonstration that a general coupled solution conserves
momentum. The nontrivial checks are the predicted relaxation and its ODE
residual. \Cref{fig:two_way_response} also shows the kinetic-energy decay,
consistent with the sign in \cref{eq:drag_dissipation}.

\ResultFigure{10_two_way_response}{V02 uniform two-way relaxation.
The two velocities approach a common value while kinetic energy decreases.
The total-momentum line is exact by the imposed output relation. Velocity
agreement and the independently sampled ODE residual assess whether the
learned temporal response satisfies the governing problem.}
{fig:two_way_response}

\section{Kernel and carrier-flow verification}
\subsection{Normalization, overlap, and parcel splitting}
The bounded-domain normalization test V03 probes periodic-seam and wall
locations. Its maximum recorded discrepancy is approximately
$8.44\times10^{-15}$ across the checks included in that test. The displayed
integrals in \cref{fig:kernel_checks} differ from one by only a few units
of machine precision. The larger summary discrepancy also includes the
constant-field sampling check and should not be assigned to an individual
integral in the figure.

V04 compares a three-dimensional product quadrature with an independent
free-space Gaussian overlap array. The relative difference is approximately
$1.29\times10^{-15}$. The CSV column retained from the numerical workflow
uses a generic prediction name; here it is correctly identified as a
quadrature result rather than a learned PINN field. This is a kernel-integral
test and does not establish a three-dimensional channel solution.

V09 splits a fixed population of 10,000 particles into 10, 50, and 100
coincident parcels. The maximum force-density difference is approximately
$2.66\times10^{-15}$. Because the split centers coincide and total mass
is unchanged, this tests deterministic representation invariance. It is not
a convergence study of independently dispersed parcel populations.

\ResultFigure{12_kernel_verification}{Gaussian transfer checks. Panel (a)
shows V04's independent three-dimensional overlap integral and product
quadrature. Panel (b) shows the stored V03 normalization values divided by
machine precision. The kernels are verified on their declared quadrature;
neither panel is a flow-field prediction.}{fig:kernel_checks}

\subsection{Unsteady incompressible flow}
V10 considers the periodic two-dimensional Taylor--Green vortex with
viscosity $\nu=0.01$ over $0\leq t\leq5$. The exact velocity is
\begin{equation}
 u=\sin x\cos y\,e^{-2\nu t},\qquad
 v=-\cos x\sin y\,e^{-2\nu t},
 \label{eq:tg_velocity}
\end{equation}
and the gauge-aligned pressure is
\begin{equation}
 p=\frac{1}{4}(\cos2x+\cos2y-2)e^{-4\nu t}.
 \label{eq:tg_pressure}
\end{equation}
The initial field is imposed, while the analytical temporal decay is held
out of training. Independent sample-relative errors are \TGU\% for $u$,
\TGV\% for $v$, and \TGP\% for pressure. The independent residual RMS is
approximately $7.38\times10^{-4}$. These results support the unsteady
carrier implementation for this smooth periodic test.

\ResultFigure{11_taylor_green}{V10 Taylor--Green verification on independent
space--time samples. Each panel compares the predicted quantity with its
analytical value, using the same pressure gauge. The test contains no
parcels and establishes carrier-flow verification rather than coupled
particle-transport validation.}{fig:taylor_green}

\begin{table}[htbp]\centering\small
\caption{Accepted solution-observable errors in the reduced Gaussian-model
verification problems. Residual criteria are checked separately.}
\label{tab:gaussian_errors}
\begin{tabularx}{\textwidth}{@{}lYr@{}}\toprule
Test & Observable & Relative error\\\midrule
\EvidenceRows{gaussian_errors}
\bottomrule\end{tabularx}
\end{table}

\section{Status of the coupled straight-channel calculations}
The evaluated coupled channel uses one training seed and four combinations
of Stokes number and mass loading. The training schedule is 1,200 Adam
steps followed by a 200-iteration L-BFGS budget; the transfer quadrature
is $16\times16$. All four calculations produce finite outputs, but none
meets the complete residual screen. \Cref{tab:coupled_scope} reports the
momentum and parcel-acceleration RMS values that bound their interpretation.

\begin{table}[htbp]\centering\small
\caption{Coupled straight-channel screening. Every required held-out
residual must be below 0.03; the momentum residual exceeds that bound in
all four cases. These are diagnostic calculations, not accepted transport
predictions.}\label{tab:coupled_scope}
\begin{tabularx}{\textwidth}{@{}rrrrY@{}}\toprule
$\St$ & $\beta$ & Momentum RMS & Parcel RMS & Complete screen\\\midrule
\EvidenceRows{coupled_scope}
\bottomrule\end{tabularx}
\end{table}

The lowest listed momentum RMS is $0.0407$, still above the declared bound.
The full screen also requires all centers to remain inside the channel,
a trained finite checkpoint without training error, local volume fraction
below 0.01, and slip Reynolds number no greater than one. The observed
maximum slip Reynolds numbers range from approximately $0.234$ to $0.642$.
These satisfy that loose screening bound but do not establish the
asymptotically small-slip regime required for accurate Stokes drag.
Passing the reduced tests therefore does not justify selecting a visually
smooth coupled-channel field as a validated result. I exclude these
predictions from quantitative claims about loading effects, sedimentation,
or deposition. The physical model also samples the disturbed fluid without
an undisturbed-velocity correction, so successful future residual screening
would still require independent assessment of the drag closure.

All ten component tests in the evaluated record pass their individual
criteria. Tests V05--V08 concern impact algebra and a frozen-rate sampling
primitive, which are outside the collision-free channel study and are not
used to validate its dynamics. The literature catalogue likewise supplies
no matched physical-validation result for the Gaussian channel. These
scope limits preserve the distinction between tested components and an
accepted coupled prediction.

\section{Chapter conclusions}
The Gaussian formulation establishes exact transfer-quadrature identities
and successful reduced tests for relaxation, kernel normalization, and
unsteady carrier flow. The strongest supported results are these
independently evaluated components. Further optimization, transfer-grid
and time-sampling refinement, multiple seeds, and a matched independent
coupled reference are required before interpreting the fully coupled
channel quantitatively. The next chapter therefore formulates a controlled
constricted-channel extension rather than extending unsupported transport
claims.
